% !TeX root = courtade-kumar.tex

We prove the properties of the comparison function used in
\autoref{sec:comparison} and the estimate used in
\autoref{lem:comparison-correlation}.
All endpoint expressions are understood by continuity.

For the following calculations, let us write
\begin{equation*}
  f(t)=\kldivb{(1+t)/2}{1/2}
  =\sum_{k\ge1}\frac{t^{2k}}{2k(2k-1)\ln2},
  \qquad -1\le t\le1.
\end{equation*}
Thus $f$ is even, convex, and strictly increasing on $[0,1]$.
Its nonnegative coefficients sum to one, so $f(t)\le t^2$.

Let $A,B$ be unbiased bits, with absolute correlation $\rho$.
Let $C$ be a nonconstant bit such that $A-B-C$ is a Markov chain,
and write
\begin{equation*}
  V=\tsent{C},\qquad
  d=\E[(2B-1)(2C-1)].
\end{equation*}
We define $F_\rho$ on $[0,1]$ by
\begin{equation*}
  F_\rho(f(t))=f(\rho t),\qquad 0\le t\le1.
\end{equation*}

\begin{lemma}
\label{lem:comparison-calculus}
The function $F_\rho$ is increasing and concave, and
\begin{equation*}
  \muti{B}{C}\le V,
  \qquad
  \muti{A}{C}\le V F_\rho\left(\frac{\muti{B}{C}}V\right).
\end{equation*}
Moreover,
\begin{equation*}
  |d|\le V,
  \qquad \muti{B}{C}\ge V f(d/V).
\end{equation*}
\end{lemma}
\begin{proof}
Let us write $x=\E[2C-1]$, so $|x|+|d|\le1$ and
$V=1-x^2$. Bayes' rule gives
\begin{align*}
  \muti{B}{C}
  &=\frac{1+x}{2}f\left(\frac{d}{1+x}\right)
  +\frac{1-x}{2}f\left(\frac{d}{1-x}\right)
  =:\ell(x,d),\\
  \muti{A}{C}&=\ell(x,\rho d).
\end{align*}
Relabeling $A$ if necessary makes its correlation with $B$
nonnegative and does not change either information.
Since $f(t)\le t^2$,
\begin{equation*}
  \ell(x,d)\le\frac{d^2}{V}\le V,
\end{equation*}
where $|d|\le1-|x|\le V$ proves the second inequality.

The divergence series also gives
\begin{equation*}
  \frac{\ell(x,d)}V
  =\sum_{k\ge1}
  \frac{W_k(x)(d/V)^{2k}}{2k(2k-1)\ln2},
  \qquad
  W_k(x)=\frac{(1+x)^{2k-1}+(1-x)^{2k-1}}2.
\end{equation*}
Here $W_1=1$ and $W_k\ge1$. This proves
$\ell(x,d)\ge Vf(d/V)$, as required for the lower bound.

Let us choose $\tau\in[0,1]$ so that $f(\tau)=\ell(x,d)/V$.
Cauchy--Schwarz gives $W_k^2\le W_{k-1}W_{k+1}$ for $k\ge2$.
Together with $W_1=1$ and $W_k\ge1$, this implies that
$W_k^{1/k}$ is nondecreasing. Indeed, successive differences of
$\ln W_k$ are nondecreasing, so
\begin{equation*}
  \ln W_{k+1}-\ln W_k\ge\frac{\ln W_k}{k-1}
  \qquad(k\ge2).
\end{equation*}
Consequently the terms
\begin{equation*}
  \frac{W_k(x)(d/V)^{2k}-\tau^{2k}}{2k(2k-1)\ln2}
\end{equation*}
change sign at most once, from negative to positive.
Their absolutely convergent sum is zero by the choice of $\tau$.
Multiplication by the decreasing weights $\rho^{2k}$ therefore
gives a nonpositive sum: subtracting a weight at the sign change
makes each summand nonpositive.
Thus
\begin{equation*}
  \ell(x,\rho d)\le Vf(\rho\tau)
  =VF_\rho\bigl(\ell(x,d)/V\bigr).
\end{equation*}
The cases with no sign change, and $\rho=0,1$, follow directly.

It remains to prove concavity. For $y=f(t)$ with $0<t<1$,
\begin{equation*}
  F_\rho'(y)=\rho^2
  \frac{\sum_{j\ge0}\rho^{2j}t^{2j}/(2j+1)}
  {\sum_{j\ge0}t^{2j}/(2j+1)}\ge0.
\end{equation*}
This ratio decreases with $t$: differentiating and pairing
terms indexed by $j,k$ gives factors
$(j-k)(\rho^{2j}-\rho^{2k})\le0$.
Since $y$ increases with $t$, the derivative decreases with $y$.
Continuity gives the assertion on $[0,1]$.
\end{proof}

Finally, the lower pointwise comparison averages in the same way.
In the setup of \autoref{sec:comparison}, let us write
$v=\tsent{C\emid U}>0$ and
$d=\E[(2B-1)(2C-1)]$.
For each value $u$ of $U$, let us write $V_u$ and $d_u$ for the
corresponding conditional entropy and correlation.
Independence of $B$ and $U$ gives $\E_{u\sim U}d_u=d$.
Since $|d_u|\le V_u$, we have $|d|\le v$.
Jensen's inequality for the convex function $f$, with weights
$V_u/v$, gives
\begin{equation*}
  \muti{B}{C\emid U}
  \ge\E_{u\sim U}\bigl[V_u f(d_u/V_u)\bigr]
  \ge v f(d/v).
\end{equation*}
Values with $V_u=0$ have $d_u=0$ and contribute zero.
